%%%%%%%%%%%%%%%%%%%%%%%%%
% Information about your work
%%%%%%%%%%%%%%%%%%%%%%%%%
%
% LaTeX uses the following information about your work to create the "front matter" of
% your document.

\title{
An Empirical Meta-Evaluation of Language Model Evaluation Methods for Systems Engineering Using Distractor Sensitivity and Consensus Judging 
}

% Student info
\author{Ryan A Bell}
\rank{Civilian, Department of the Navy}    %\rank{Civilian} % if you don't have a rank
\degree{Doctor of Philosophy in Systems Engineering}
\degreeabbreviation{PhD}   % Should be MS, MBA or MA
\prevdegrees{BSEE, MSEE} % previous degree

% Department info
\department{Department of Systems Engineering}
\thesisadvisor{Dr. Raymond Madachy}
\memberOne{Dr. Raymond Madachy}
\memberTwo{Dr. Raymond Madachy}
\memberThree{Dr. John Green}
\memberThreeTitleLine{Department of Systems Engineering}
\memberFour{Dr. Douglas Van Bossuyt}
\memberFourTitleLine{Department of Systems Engineering}
\memberFive{Dr. Ronald Giachetti}
\memberFiveTitleLine{Department of Systems Engineering}
\memberSix{Dr. Mathias Kolsch}
\memberSixTitleLine{Department of Computer Science}

\departmentchair{Oleg A. Yakimenko}

% The date you are graduating:
\degreedate{March 2026}

% Use same distribution statement selected on your thesis dashboard.
\distribution{Approved for public release. Distribution is unlimited.}

\acresetall
% Paste abstract from your thesis dashboard. New paragraphs start after an empty line.
\abstract{The prominence of Large Language Models (LLMs) has opened new avenues in systems engineering, enhancing capabilities in requirements analysis, design synthesis, and decision support. Benchmarking these models effectively remains a challenge due to the complexity and specificity of the domain. This paper explores the comparative effectiveness of Multiple Choice Questions (MCQs) and Textual Responses as methods for evaluating LLM performance in systems engineering contexts. Utilizing a tailored Systems Engineering (SE) language model benchmark tailored for domain specific tasks, we assess state-of-the-art LLMs across both evaluation modalities.
To create a Textual Responses dataset from our original Multiple Choice Questions (MCQs) and ensure an "apples-to-apples" comparison, we employ an LLM-assisted approach to convert structured MCQs into open-ended questions. This method involves having an LLM generate textual questions that retain the core concepts of each MCQ but are rephrased to elicit a detailed, free-text response instead of a multiple-choice selection. The LLM-in-the-loop process allows us to automate the creation of an extensive and varied textual dataset while maintaining fidelity to the original concept in question. To compare the responses, we leverage cosine similarity metrics to quantify the semantic alignment between the models' answers and the reference solutions. LLM as a Judge for grading alignment will also be used. Both approaches enable measuring the degree of conceptual overlap between response and reference solution in a high-dimensional vector space in an attempt to reflect the models' understanding of complex systems engineering concepts. Analysis of each evaluation technique will be discussed.
}

% If CUI (Controlled Unclassified Information), swap which line is commented-out (the % sign) for your drafts.
\securitybanner{}
% \securitybanner{CUI} % remove before requesting FINAL REVIEW with thesis processor, who will insert final banners.
%
% Do NOT upload CUI to public websites like Overleaf.com.
%
% If you are interested in getting a SECURE Overleaf site approved for NPS CUI work, visit https://nps.edu/web/thesisprocessing/latex-cui.
%
% Mandatory fields for the SF298.
%
\ifnpstechreport
    \ReportType{Technical Report}
\else
    \ReportType{Dissertation}
\fi
\ReportDate{[Month Year]}       % for a thesis, graduation date
\SponsoringAgency{N/A}          % really, for technical reports
\DatesCovered{MM-DD-YYYY to MM-DD-YYYY}
\ReportClassification{Unclassified}
\AbstractClassification{Unclassified}
\PageClassification{Unclassified}
%
% Optional fields for the SF298.
%
\RPTpreparedFor{}
\ContractNumber{}
\GrantNumber{}
\ProgramElementNumber{}
\TaskNumber{}
\WorkUnitNumber{}
\POReportNumber{}
\Acronyms{}
\SMReportNumber{}
\SubjectTerms{TBD}
\ResponsiblePerson{}
\RPTelephone{}
\SignatureOne{}
\SignatureTwo{}
\SupplementaryNotes{The views expressed in this document are those of the author and do not reflect the official policy or position of the Department of Defense or the U.S. Government.
}

% Optional. Prevents footnotes from being reset at each chapter
% Comment this out to have them reset with each chapter.
\makeatletter
\@removefromreset{footnote}{chapter}
\makeatother

% Optional. Adds pdf metadata and links.
% This should be right before the \begin{document}, to be the
% last package / macros defined. (Hyper-ref is fragile,
% needs to be last, and has known conflicts with other packages.)
% Comment out if you have build problems building with hyperref
%\NPShyperref


%%%%%%%%%%%%%%%%%%%%%%%%%
% Your thesis begins here.
%%%%%%%%%%%%%%%%%%%%%%%%%

\begin{document}
%% IEEE format only and Mechanical and Aerospace Engineering (MAE) Department only:
%% Uncomment the next line.  All others, leave commented out (with % sign)
%\bstctlcite{forced-et-al-off}
\NPScover                  % Cover page
\NPSsftne                  % SF298 form
%\NPSsignature             % Tech Report page (iii): signature page
\ifnpstechreport
    \else
    \NPSthesistitle            % Thesis page (iii): title page
    \fi
\NPSabstractpage           % Abstract Page
\NPSfrontmatter            % NPS front matter follows

\ifnpstechreport
\else
\DraftwatermarkOptions{stamp=false} % Starting at this point in the document, this commands the draftwatermark package to stop printing.
\fi

% This changes the chaptermark and includes the various tables
% It must be here.  Comment this line out for articles that do not contain chapters.
\renewcommand{\chaptermark}[1]{\markboth{\MakeUppercase{\chaptername}\ \thechapter.\ #1}{}}

%
% If you don't need one of these, comment it out.
%
%\setcounter{page}{7} % -- if you need to change the page number for the TOC do it here
\NPStableOfContents
\NPSlistOfFigures
\NPSlistOfTables
\NPSlistOfAcronymsFromFile{acronyms}

%
% Put (optional) Executive Summary here.
% New paragraphs start after an empty line. See exec_sum_with_refs.tex under the examples directory for a demo.
% If not using, comment-out or delete all of these lines.
%
\acresetall
%\NPSexecsummary{%
%This is different from your abstract. An executive summary is required or recommended for the departments listed \href{https://nps.edu/documents/105790666/106471207/Executive_Summary_Guidance.pdf}{\underline{here}}. Most executive summaries range from 2--5 pages.

%An executive summary is a highly condensed version of your thesis that should be able to stand alone, independent of your thesis. An executive summary should summarize your purpose, methods, results, conclusions, and recommendations to allow someone who can read ONLY that document to walk away with a solid understanding of the research.

%If you include figures or tables in your executive summary, do not use the caption commands for the titles. Instead, manually type in the titles. This will keep these titles out of your thesis's Lists of Figures and List of Tables, and allow the figure and table numbering to start at “1” in the thesis body, as required.

% How to have References in your Executive Summary:
%If you include citations in the Executive Summary, it may make more sense to type your entries in manually, since most executive summaries contain only a handful of sources. To add a reference list programmatically, see exec\_sum\_with\_refs.tex in the ``additional\_resources'' node of this template for the code.
% Instead of having this sample Executive Summary content here, put your executive summary in the
% file additional_resources/exec_sum_with_refs.tex and use this next line to link it here:
%\input{additional_resources/exec_sum_with_refs}
%}

%
% Put (optional) acknowledgements here.  
% New paragraphs start after an empty line.
% If not using, comment-out or delete all of these lines.
%
\acresetall
\NPSacknowledgements{I would like to express my deepest gratitude to Dr. Raymond Madachy for his mentorship, generosity with his time, and steady support throughout this work. He helped me grow into academic life by encouraging exploration, challenging assumptions, and showing what it means to collaborate with intellectual honesty and respect. His enthusiasm for new topics and his ability to bridge emerging AI concepts with established systems engineering practice directly shaped this dissertation’s scope and direction. He also demonstrated a model of professorship I hope to emulate: rigorous without being rigid, practical without sacrificing depth, and comfortable navigating complexity where the right answer is rarely black-and-white.

I am grateful to Dr. Rajendra Singh, who challenged me long ago to choose between the technical route (a PhD) and the management route (an MBA), and who opened my eyes to how often we underestimate what’s possible by not thinking big enough. 

I am also sincerely thankful to my dissertation committee: Dr. Ronald Giachetti, Dr. John (Mike) Green, Dr. Mathias Kolsch, and Dr. Douglas Van Bossuyt for their collaborative spirit, thoughtful questions, and willingness to provide alternative viewpoints that helped refine and strengthen these ideas. Their guidance ensured this research stayed grounded, useful, and oriented toward benefiting the broader systems engineering community.

In the same spirit, I am especially grateful to Ryan Longshore for sharing in a deeply curious approach to exploring new concepts through late-night brainstorms, candid debate, and the kind of devil’s advocate conversations that sharpen ideas into something clearer, more defensible, and more worth sharing.

To my wife, thank you for your unwavering support and for the courage and conviction with which you chase what’s next; your spirit inspires me and your reminders to keep reaching higher keep me moving forward, especially when the path feels steep, uncertain, or inconvenient.

To my family, thank you for your timeless support and for instilling in me an empathetic understanding of serving something larger than myself: service to country, commitment to community, and a responsibility to contribute to the greater good.

Finally, I am grateful to the many colleagues, mentors, and advocates across NIWC, NPS, and my personal circles who challenged my assumptions, broadened my perspective, and helped shape the viewpoints I carry today. This work reflects not only my effort, but the collective influence of everyone who invested their time and insight along the way.

}

% Start layout for the NPS body
\NPSbody
